\section{Conceptos fundamentales del análisis estadístico}
    \subsection{Distribución de frecuencias}
        Una distribución de frecuencias organiza los datos en clases o intervalos y
        registra cuántas observaciones caen en cada una \cite{triola2018estadistica},
        \cite{walpole2012probabilidad}, \cite{sturges1926choice}, su lectura se apoya
        en cuatro componentes: la frecuencia absoluta ($f_{i}$) que cuenta las
        observaciones de la clase, la frecuencia acumulada ($F_{i}$) que suma las
        absolutas hasta esa clase, la frecuencia relativa ($h_{i}$) la cual expresa la
        proporción que la clase representa sobre el total y la relativa acumulada
        ($H_{i}$) que indica qué fracción de los datos queda por debajo del límite
        superior de la clase \cite{triola2018estadistica},
        \cite{walpole2012probabilidad}, adicionalmente, cada clase se representa
        mediante su marca de clase, que corresponde al punto medio del intervalo
        \cite{triola2018estadistica}, \cite{walpole2012probabilidad},
        \cite{sturges1926choice}, \cite{few2012show}, \cite{castro2026visualizacion}.
        Para este informe, el número de clases ($k$) se obtiene a partir de la regla de
        Sturges, que se define en la ecuación (\ref{sturges}).

        \begin{equation}
            k = 1 + 3{,}322 \, \log_{10}(n)
            \label{sturges}
        \end{equation}

        Con $n = 120$ observaciones se hallan cerca de 8 clases de amplitud constante,
        con las cuales se construye la tabla de frecuencias del consumo energético. La
        distribución de frecuencias tiene además tres representaciones gráficas directas
        que se elaboran en este informe que son:

        \vspace{10pt}

        \begin{itemize}
            \item \textbf{Histograma}: muestra barras contiguas cuya altura corresponde
            a la frecuencia absoluta \cite{castro2026visualizacion}.

            \item \textbf{Polígono de frecuencias}: une con una línea las marcas de
            clase \cite{castro2026visualizacion}.

            \item \textbf{Ojiva}: traza la curva de la frecuencia acumulada sobre la
            cual puede leerse la mediana en el 50\% de los datos
            \cite{castro2026visualizacion}.
        \end{itemize}

    \subsection{Medidas de tendencia central}
        Las medidas de tendencia central localizan el valor alrededor del cual se agrupa
        la distribución, algunas de ellas son: la media aritmética, que corresponde al
        promedio de las observaciones y aprovecha toda la información disponible, no
        obstante, es sensible a los valores extremos, la mediana que obtiene el valor
        que divide los datos ordenados en dos mitades iguales y se mantiene robusta
        frente a esos extremos, por último, la moda es el valor que más se repite
        \cite{triola2018estadistica}, \cite{walpole2012probabilidad},
        \cite{castro2026visualizacion}.

        \vspace{10pt}

        Cuando la variable es continua, la moda puntual carece de sentido práctico, por
        lo que se reporta la clase modal, es decir, la de mayor frecuencia absoluta,
        junto con la moda interpolada dentro de ella mediante la expresión de la
        ecuación (\ref{moda_interpolada}).

        \begin{equation}
            M_{o} = L + \frac{d_{1}}{d_{1} + d_{2}} \, w
            \label{moda_interpolada}
        \end{equation}

        Donde $L$ es el límite inferior de la clase modal, $w$ su amplitud y $d_{1}$ y
        $d_{2}$ las diferencias de frecuencia con las clases adyacentes
        \cite{walpole2012probabilidad}. Si la variable es nominal, como ocurre con el
        sector, la moda es simplemente la categoría más frecuente.

        \vspace{10pt}

        La relación entre las tres medidas funciona como un diagnóstico de la forma de
        la distribución, ya que cuando la media, la mediana y la moda coinciden, la
        distribución es aproximadamente simétrica, mientras que, si la media supera a la
        mediana y esta a su vez a la moda, la distribución presenta asimetría positiva o
        cola derecha \cite{triola2018estadistica}, \cite{castro2026visualizacion}.

    \subsection{Medidas de dispersión}
        Las medidas de dispersión cuantifican qué tanto se alejan las observaciones del
        centro de la distribución. La más elemental es el rango, que corresponde a la
        diferencia entre el valor máximo y el mínimo de la muestra y depende únicamente
        de los extremos \cite{triola2018estadistica}, \cite{walpole2012probabilidad},
        \cite{castro2026visualizacion}, como lo expresa la ecuación (\ref{rango}).

        \begin{equation}
            R = x_{\max} - x_{\min}
            \label{rango}
        \end{equation}

        La varianza muestral promedia las desviaciones al cuadrado de cada observación
        respecto de la media aritmética y divide entre $n - 1$ para obtener un estimador
        sin sesgo \cite{triola2018estadistica}, \cite{walpole2012probabilidad},
        \cite{castro2026visualizacion}, de acuerdo con la expresión matemática de la
        ecuación (\ref{varianza}).

        \begin{equation}
            s^{2} = \frac{1}{n - 1} \sum_{i=1}^{n} \left( x_{i} - \bar{x} \right)^{2}
            \label{varianza}
        \end{equation}

        Donde $x_{i}$ es cada observación, $\bar{x}$ la media aritmética y $n$ el tamaño
        de la muestra. La desviación estándar está formulada en la ecuación
        (\ref{desviacion}) y corresponde a la raíz cuadrada de la varianza, a partir de
        esa ecuación se recuperan las unidades originales de la variable, lo que
        facilita su interpretación directa \cite{triola2018estadistica},
        \cite{walpole2012probabilidad}, \cite{castro2026visualizacion}.

        \begin{equation}
            s = \sqrt{s^{2}} =
            \sqrt{\frac{1}{n - 1} \sum_{i=1}^{n} \left( x_{i} - \bar{x} \right)^{2}}
            \label{desviacion}
        \end{equation}

        Adicionalmente, el coeficiente de variación es presentado en la ecuación
        (\ref{cv}) y expresa la desviación estándar como porcentaje de la media 
        permitiendo comparar la variabilidad de grupos con escalas distintas
        \cite{walpole2012probabilidad}.

        \begin{equation}
            CV = \frac{s}{\bar{x}} \times 100\,\text{\%}
            \label{cv}
        \end{equation}

        Por último, el rango intercuartílico de la ecuación (\ref{iqr}) corresponde a 
        la diferencia entre el tercer y el primer cuartil, mide la amplitud del 50\%
        central de los datos y constituye la base del diagrama de caja
        \cite{walpole2012probabilidad}, matemáticamente se define como:

        \begin{equation}
            IQR = Q_{3} - Q_{1}
            \label{iqr}
        \end{equation}
